\chapter{Deviated Septum}

\section*{Definition and Overview}
The nasal septum is the wall of cartilage and bone that divides the nose into two passages. A deviated septum occurs when this wall is bent or crooked, so that one nasal passage is narrower than the other. Many people have a slight septal deviation that causes no problems, but a more significant deviation can cause a blocked nose, difficulty breathing, snoring, and recurrent sinus problems. A deviated septum can be present from birth, develop during childhood growth, or result from injury to the nose. This chapter covers the causes, symptoms, and treatment options.

\section*{Epidemiology}
A deviated septum is very common: studies suggest that up to 80\% of people have some degree of septal deviation, although only a minority have symptoms severe enough to seek treatment. It is equally common in men and women. Nasal injuries are a common cause, especially in adolescents and young adults from sport, falls, and fights. Deviated septum is one of the most common reasons for referral to an ear, nose, and throat (ENT) surgeon, and septoplasty is a frequently performed operation.

\section*{Aetiology and Pathophysiology}
\begin{itemize}
    \item \textbf{Congenital factors:} the septum can be crooked from birth due to the way the nose formed in the womb
    \item \textbf{Growth:} uneven growth of the septum and surrounding facial structures during childhood can cause deviation
    \item \textbf{Trauma:} injury to the nose, including during birth or from falls, sport, and assaults, can bend or displace the septum
    \item \textbf{Effects:} a deviated septum narrows one airway, disrupting airflow and creating turbulence that can dry out the lining
    \item \textbf{Associated structures:} enlargement of the turbinates (bony shelves in the nose) on the opposite side is common
    \item \textbf{Impact:} reduced airflow, nasal blockage, and impaired drainage of the sinuses
\end{itemize}

\section*{Clinical Features}
\begin{itemize}
    \item A persistently blocked or congested nose, often worse on one side
    \item Difficulty breathing through the nose, especially during exercise or sleep
    \item Nosebleeds, which can result from the dry, turbulent airflow
    \item Snoring and sleep-disordered breathing
    \item Recurrent sinus infections, because sinus drainage is impaired
    \item Facial pain or pressure in some people
    \item The visible appearance of a bent nose after injury
    \item Many people with mild deviation have no symptoms at all
\end{itemize}

\section*{Differential Diagnosis}
\begin{itemize}
    \item Allergic rhinitis, which causes a blocked, runny nose without a structural problem
    \item Nasal polyps, which can obstruct the nasal passages
    \item Enlarged turbinates causing blockage
    \item Chronic sinusitis, which can coexist with a deviated septum
    \item Foreign body in the nose, especially in children
    \item Nasal tumours, which are rare but serious
    \item Sleep apnoea, which can cause snoring and daytime sleepiness
\end{itemize}

\section*{When to Seek Medical Care}
See a doctor if a blocked nose is persistent, affects breathing, sleep, or daily life, or is associated with recurrent sinus infections or nosebleeds. A GP can assess the nose and refer to an ENT specialist if needed. Anyone with nasal obstruction after injury, or with one-sided symptoms, should also seek review.

\textbf{Call 000 (or local emergency services) for:}
\begin{itemize}
    \item A nosebleed that does not stop with firm pressure for 20 minutes
    \item Difficulty breathing through both nostrils or shortness of breath
    \item Significant facial injury with deformity, vision changes, or bleeding that will not stop
    \item A suspected foreign body causing breathing difficulty in a child
\end{itemize}

\section*{Diagnosis and Investigations}
\begin{itemize}
    \item \textbf{History:} symptoms, injuries, and their impact on breathing and sleep
    \item \textbf{Nasal examination:} inspection of the nasal passages with a speculum
    \item \textbf{Nasendoscopy:} examination with a small flexible telescope to assess the septum, turbinates, and sinuses
    \item \textbf{Imaging:} CT scans are not usually needed but may be used to assess the sinuses or plan surgery
    \item \textbf{Sleep studies:} if sleep apnoea is suspected
    \item \textbf{Allergy testing:} where allergic rhinitis contributes to symptoms
\end{itemize}

\section*{Management}
\begin{itemize}
    \item \textbf{Observation:} mild, asymptomatic deviation needs no treatment
    \item \textbf{Medical management:} nasal corticosteroid sprays reduce swelling of the lining and turbinates, improving airflow
    \item \textbf{Decongestants and saline rinses:} short-term symptom relief and moisturising
    \item \textbf{Treatment of associated conditions:} allergy management and treatment of sinusitis
    \item \textbf{Septoplasty:} surgery to straighten the septum, performed under general or local anaesthetic
    \item \textbf{Combined procedures:} septoplasty is often combined with turbinate reduction or sinus surgery
    \item \textbf{Recovery:} nasal packing or splints, with most people returning to normal activities within a week or two
\end{itemize}

\section*{Complications}
\begin{itemize}
    \item Chronic nasal blockage and mouth breathing, which affects sleep and quality of life
    \item Recurrent sinusitis and nosebleeds
    \item Snoring and sleep-disordered breathing
    \item In children, effects on facial growth and school performance from mouth breathing
    \item Surgical complications of septoplasty, including bleeding, infection, perforation of the septum, and changes in sensation, which are uncommon
    \item Persisting symptoms if medical treatment is not tried or surgery is not appropriate
\end{itemize}

\section*{Prognosis and Outcomes}
For people with symptomatic septal deviation, treatment is very effective. Many improve with medical therapy alone, and septoplasty provides significant and lasting improvement in nasal breathing for the majority of people who need it. Symptoms such as blockage, snoring, and sinus infections usually improve substantially. Complications are uncommon, and most people are satisfied with the outcome. Even with treatment, some degree of nasal congestion can persist, particularly with allergies or other factors.

\section*{Prevention and Self-Care}
\begin{itemize}
    \item Protect the nose from injury with mouthguards and face protection in contact sport
    \item Use seatbelts and appropriate restraints to reduce facial injuries
    \item Use nasal corticosteroid sprays as prescribed for congestion
    \item Avoid picking the nose, which can cause nosebleeds
    \item Use saline rinses to keep the nose moist and clear
    \item Manage allergies with appropriate treatment
    \item Seek ENT review if symptoms persist despite medical treatment
    \item Follow the surgeon's advice carefully if you have surgery
\end{itemize}

\section*{Sources and Further Reading}
The following reputable sources provide further information for healthcare students and patients seeking reliable, current advice about a deviated septum.

\begin{itemize}
    \item Australian Society of Otolaryngology Head and Neck Surgery. \textit{Septoplasty and nasal obstruction information.} https://www.asohns.org.au/
    \item NHS. \textit{Deviated septum: symptoms and treatment.} https://www.nhs.uk/conditions/deviated-nasal-septum/
    \item Mayo Clinic. \textit{Deviated septum: diagnosis and treatment.} https://www.mayoclinic.org/diseases-conditions/deviated-septum/
    \item Healthdirect Australia. \textit{Deviated septum and nasal blockage.} https://www.healthdirect.gov.au/
    \item MedlinePlus. \textit{Deviated septum.} https://medlineplus.gov/
    \item MSD Manual. \textit{Nasal obstruction: professional overview.} https://www.msdmanuals.com/
\end{itemize}
